\section{总结与延伸}

\subsection{本讲核心要点回顾}

MIT 6.S191 第一讲从宏观到微观，系统地介绍了深度学习的基础框架：

\begin{enumerate}
\item \textbf{深度学习的定位}：AI $\supset$ ML $\supset$ DL，核心是从原始数据中自动学习特征表示
\item \textbf{感知机}：神经网络的基本单元，包含加权求和、偏置和非线性激活三个步骤
\item \textbf{激活函数}：Sigmoid、Tanh、ReLU 等非线性函数使网络能够逼近任意复杂函数
\item \textbf{深度网络}：通过堆叠多个 Dense 层构成，每层学习越来越抽象的特征
\item \textbf{损失函数}：交叉熵（分类）和 MSE（回归）分别衡量预测与真实值的差距
\item \textbf{训练过程}：梯度下降 + 反向传播 + 适当的学习率
\item \textbf{正则化}：Dropout 和 Early Stopping 防止过拟合
\end{enumerate}

\begin{table}[H]
\centering
\begin{tabular}{lp{5cm}p{5cm}}
\toprule
\textbf{概念} & \textbf{数学表达} & \textbf{直觉解释} \\
\midrule
感知机 & $\hat{y} = g(w_0 + \mathbf{X}^T\mathbf{W})$ & 加权投票 + 非线性决策 \\
Dense 层 & $z_i = w_{0,i} + \sum_j x_j w_{j,i}$ & 多个并行的感知机 \\
损失函数 & $J(\mathbf{W}) = \frac{1}{n}\sum_i \mathcal{L}$ & 预测质量的量化度量 \\
梯度下降 & $\mathbf{W} \leftarrow \mathbf{W} - \eta \nabla J$ & 沿最陡方向下山 \\
反向传播 & 链式法则递归应用 & 从输出向输入逐层传递误差 \\
Dropout & 随机置零概率 $p$ & 强制冗余，防止依赖单一特征 \\
\bottomrule
\end{tabular}
\caption{本讲核心概念速查表}
\end{table}

\subsection{从基础到前沿}

Amini 在讲座中多次强调，本讲的内容（感知机、梯度下降、反向传播）都是数十年前的发明，但它们是理解现代深度学习的必要基础。课程后续将涵盖：

\begin{itemize}
\item \textbf{序列建模}：RNN、LSTM 等处理时序数据的架构
\item \textbf{计算机视觉}：CNN 和现代视觉架构
\item \textbf{生成模型}：VAE、GAN、Diffusion Models
\item \textbf{强化学习}：让模型通过与环境交互来学习策略
\item \textbf{大语言模型}：Transformer 架构与现代 LLM
\end{itemize}

\subsection{拓展阅读}

\begin{itemize}
\item 课程官网：\url{https://introtodeeplearning.com}
\item 实验代码仓库：\url{https://github.com/MITDeepLearning/introtodeeplearning}
\item Ian Goodfellow, Yoshua Bengio, Aaron Courville, \textit{Deep Learning}, MIT Press, 2016
\item Michael Nielsen, \textit{Neural Networks and Deep Learning}, 在线免费教材
\item 3Blue1Brown, \textit{Neural Networks} 系列视频（直觉可视化）
\end{itemize}

%% --- Fin del contenido --- %%

\end{document}
